\scanpage{85}
\mspage{10}
Un cas trivial intéressant est celui où $f$ est un isomorphisme --- et $=$ celui où $X=Y$, et où $f$ fait partie d'un groupe $\unsure{\mathcal G|G}$ de $K$-isomorphismes de $X$ --- alors on peut considérer $\Pi_1(X)$ et $\Pi_1(X_{\bar K}),\Pi_1(X_K)$ comme des groupoïdes où opère $\unsure{\mathcal G|G}$, et $\widehat\Pi_1(X)$ comme un groupoïde où opère $\unsure{\mathcal G|G}\times\unsure{\mathbf{\Gamma}|\Gamma}$, $\widehat\Pi_1(\unsure{\bar X|X}_{\bar K})$ comme un groupoïde sur lequel $\unsure{\mathcal G|G}\times\Gamma$ opère. Cette fois-ci le groupe n'opère pas trivialement sur les sommets de $\Pi_1(X_K)$, l'opération de $\Gamma$ est triviale, mais pas nécessairement celle de $\unsure{\mathcal G|G}$.

Si on se donne un pt $x_0\in X_K$ comme pt base, et un ens de sommets $I\subset X_K$ contenant $x_0$, l'opération de $\Gamma$ sur $\widehat\Pi_1(X,I)$ est connue quand on connaît son opération sur $\widehat\pi_1(X,x_0)$, et si de plus on connaît son opération sur les $\widehat\pi_1(X,x_0)$-torseurs à droite
\figure{\begin{tikzpicture}\draw(0,0)--(0,.65)--(1.3,.78)--(1.3,1);\end{tikzpicture}}
\[
\tag{41}\uline{\operatorname{Hom}}_{\widehat\Pi_1(X,I)}(x_0,x_i),\qquad i\in I.
\]
ou, ce qui revient au $=$, si on connaît les
\margin{84}
